
\section{Methodology}
\label{sec:setup}

\subsection{Dataset and Control Variants}
\label{sec:dataset}

\paragraph{Question pool, taxonomy, and selection.}
We begin with 1{,}000 questions randomly sampled from the VQA v2.0 validation split
\citep{goyal2016making}, spanning \textit{yes/no}, numerical, and open-ended answer types.
% The original VQA metadata question types are grouped by their surface question form (e.g., \textit{how many}, \textit{what color}), 
We annotate each question along two orthogonal dimensions: \emph{operation} and \emph{entity} type (full taxonomy in Appendix~\ref{app:taxonomy}). 
We then diversity-select 113 questions spanning all 11 operation types and 9 entity types (full composition in Table~\ref{tab:q113_sample_composition}). % ; question space coverage in Appendix~\ref{app:tsne} Appendix~\ref{app:sample_composition},  
For group-level anchoring analyses, operation and entity types with $n \leq 5$ questions are collapsed into an \emph{Other} category (\S\ref{sec:human_patterns}). % (\textit{operations}: World Knowledge, Text Reading, Temporal, Comparison, Causality; \textit{entities}: Text, Vehicle, Place, Product); 
All groups are retained in semantic alignment analysis (\S\ref{sec:alignment_main}). %  at full granularity 
% Yes/no questions were intentionally included to capture presence-bias and count-default behaviors, but they are not evenly distributed across operation groups---\textit{Existence} is 100\% yes/no, \textit{Action} 72\%, and \textit{World Knowledge} 50\%---so SBERT-based anchor sensitivity is only interpretable for predominantly open-ended groups (\textit{Identity}, \textit{Spatial}, \textit{Attribute}, \textit{Causality}); yes/no groups are analyzed distributional in \S\ref{sec:dist_alignment}.

\paragraph{Control question variants.}
\label{sec:control}
To test how strongly blind answers depend on entity-denoting lexical anchors,
we generate three progressively abstracted variants of each question,
indexed by $v\!\in\!\{\vOrig,\vW,\vPron\}$ (referred to as \emph{variants} throughout):
\textbf{Original} — the unmodified question;
\textbf{Weaker} — object nouns generalized to a higher-level category (e.g., \emph{dog} $\rightarrow$ \emph{animal});
\textbf{Pronominalized} — content noun phrases replaced with minimal pronouns (e.g., \emph{dog} $\rightarrow$ \emph{it}) while retaining grammatical structure.
The full four-step generation schema and transformation rules are in Appendix~\ref{app:ctl}.
% When an anchor is present the referent is unambiguous and answers tend to converge; removing it widens the plausible referent space, causing answer dispersion.
% A drop in $\sHH$ similarity across variants therefore reflects this increased dispersion rather than a change in question difficulty per se.

% These variants define a controlled abstraction hierarchy that progressively
% reduces entity-specific information while preserving the operation required to
% answer the question. This allows us to test whether human and model responses
% depend primarily on lexical cues or on broader semantic and commonsense priors.
% The same ladder also functions as a pronominalization test for anchor
% dependence: if blind behavior degrades sharply when content-noun anchors are
% weakened or replaced by pronouns, then the model is relying on surface lexical
% anchors rather than stable, human-like scene priors.

\subsection{Experimental Setup}
\label{sec:exp_setup}

\paragraph{Models and Inference.}

We evaluate 20 checkpoints spanning scales of 0.6B--32B across three groups of models: \textbf{VLMs}, \textbf{backbone decoders (LM)}, and \textbf{standalone LLMs (LLM)}. 
Qwen3 LLM is evaluated both with and without Chain-of-Thought reasoning traces (as \textit{Think} model). 
SA-LLMs are text-only; several share a family with our VLMs in use. 
Appendix Table~\ref{tab:model_inventory} lists out all models we used for the analysis.

All evaluations are zero-shot at temperature~0 with a maximum of 512 output tokens, run on NVIDIA TITAN RTX GPU. 
For VLMs, they require inputs of images so we provide a blank 224$\times$224 image, where backbone decoders run the same weights text-only with image-encoder bypassed.  %  (Qwen3 \citep{yang2025qwen3}, Qwen2.5-Instruct, Vicuna)

We use 7-8B scale models for the main-text alignment tests. Two LLMs -- Phi-3.5-mini (3.8B) and Mistral-7B -- were excluded from our analysis due to extreme verbose answers (Phi: 14--16, Mistral: 28--30 on average). % (Appendix Table~\ref{tab:comprehensive_inst_blind}; Table~\ref{tab:abstention_by_qtype} , Humans: 1.1 words 

%Main-text alignment tests use the matched 7--8B subset; full-scale results are in Appendix~\ref{app:scale_variant}.
%  (Qwen3-VL \citep{bai2025qwen3vl}, LLaVA \citep{liu2023visual,liu2024llavanext}, InternVL3.5 \citep{wang2025internvl35})

\paragraph{Human Data Collection.}\label{sec:human_data}

40 participants (online university forums, IRB-approved) answered 339 (=113 original questions $\times$ 3 variants) questions across all three control variants. %  under the same blind-awareness instruction as the models 
Full protocol details and statistics are in Appendix~\ref{app:human_study}. % ---\textit{``No image provided---imagine a typical scene and answer accordingly''}. % ---yielding 13,560 total responses (339 questions $\times$ 40 humans).

\paragraph{Scenario Instruction}
For both models and humans, we prompt \textit{``No image provided---imagine a typical scene and answer accordingly,''}. 
% ; we report instruction effect on abstention as $\dinst \equiv \mathrm{Abst}(\text{w/}) - \mathrm{Abst}(\text{w/o})$ 
Instruction effects on abstention and alignment are reported in Appendix~\ref{app:abstention}. % This generally reduces refusals (full analysis in Appendix~\ref{app:abstention}).  
% Unless otherwise stated, all reported values use the with-instruction condition.%; the no-instruction baseline appears as $\Delta$ where shown.

\subsection{Metrics and Evaluation Protocols}
\label{sec:eval_protocols}

\textbf{Preprocessing.} Before scoring, all answers are normalized with the standard VQA pipeline
\citep{antol2015vqa}, and stripped of generic wrappers so human and model
responses are compared on the same cleaned answer string.
We then use regex classifier to detect abstention responses and exclude them for model-human alignment analysis. Details are in Appendix~\ref{app:preprocessing}. %  (see Appendix~\ref{app:abstention_classifier} for exclusion criteria)
%  before all distributional and semantic alignment metrics  % ; uncertain but still committed answers (e.g., \textit{``maybe''}, \textit{``perhaps''}, \textit{``unsure''}) are retained 
% We additionally report abstention ($\mathrm{Abst}(\%)$) as a human-alignment signal.
In our study, participants were instructed to provide an answer in all cases, yielding a near-zero human abstention rate ($<\!0.4\%$). % ; we assume that without such instruction, humans would naturally abstain on most questions when no image is available.
% We therefore treat blind-condition abstention (w/o) as alignment with this hypothetical uninstructed human behavior---higher is better---and instruction-condition abstention (w/) as alignment with observed human behavior---lower is better, approaching the human near-zero rate.
We report the abstention rate of models in Appendix~\ref{app:abstention}. % and instruction-effect and~\ref{app:inst}  Table~\ref{tab:alignment_ranking_7b_ticks} and 

\textbf{Distributional alignment.} For categorical (\textit{yes/no} and \textit{count}) questions that are naturally narrow in answer space, we measure alignment using Jensen--Shannon divergence ($\djs$) and a modal-match rate (\textit{Mode}).

Let $\hat{P}_{H}^{s}$ and $\hat{P}_{M}^{s}$ denote the empirical answer distributions of humans and model $M$, pooled over all questions $q\!\in\!\mathcal{Q}_{s}$ and control variants within answer-type subgroup $s\!\in\!S=\{\text{yes/no},\,\text{count}\}$; these estimate the subgroup-averaged priors $\mathbb{E}_{q}[P_{H}(a\mid q)]$ and $\mathbb{E}_{q}[P_{M}(a\mid q)]$ (or $\mathbb{E}_{q}[P_{M}(a\mid q,v_{0})]$ for VLMs).
Divergence per subgroup is
\begin{equation}
  \djs[M,s] = \djs\!\bigl(\hat{P}_{H}^{s} \,\|\, \hat{P}_{M}^{s}\bigr),
  \label{eq:js}
\end{equation}
averaged over answer-type subgroups $s$, where $\mathcal{Q}_s$ is the set of questions of type $s$.
% This captures group-level biases rather than per-question agreement, for which modal-match rate is the more direct measure.

Mode complements JS by revealing per-question majority alignment. 
For each (question, variant) pair, Mode checks whether the model's answer matches the human majority answer for that specific question, then averages over all pairs under the same sub-group.
% Unlike yes-rate---the aggregate proportion of yes answers---Mode is sensitive to \textit{which} questions the model gets right: a model with a matching aggregate yes-rate can still score low on Mode if it answers yes on the questions where humans predominantly say no, and vice versa.
% Mode thus captures per-question majority alignment rather than aggregate distributional tendency. % : a model can match the aggregate yes-rate while scoring poorly on Mode if it answers yes on the wrong questions ; conversely, a no-biased model can achieve high Mode by correctly identifying which specific questions humans predominantly answer no to. 

As a calibration reference, we report a leave-one-out (LOO) baseline: each participant's pooled distribution within $s$ is compared against the remaining participants', then averaged over participants. The resulting 95\% bootstrap CI serves as a target range---a model within this range matches human distributional behaviour.
We refer to this interval as the \textit{human range} throughout.
For $\djs$, a model below the lower bound is closer to the pooled human distribution than our human set are, indicating an majority-biased response pattern rather than stronger human-likeness.% ; above the upper bound indicates weaker alignment to the pooled human distribution.

\textbf{Semantic alignment.} For open-ended questions we compute pairwise Sentence-BERT (\citealt{reimers2019sentence}) cosine similarity using the \texttt{all-mpnet-base-v2} encoder \citep{song2020mpnet}, denoted as SBERT ($S$).

Let $\phi(\cdot)$ denote the encoder and $H$ the set of participants.
Here, $q$ indexes a single question instance, i.e., one original question under one specific control variant.
For model $m$ and question instance $q$:
\begin{equation}
  \sHM[m,q] = \frac{1}{|H|}\sum_{h\in H}\cos\!\bigl(\phi(\hat{a}_{m}(q)),\,\phi(a_{h}(q))\bigr),
  \label{eq:sbert}
\end{equation}
where $\hat{a}_{m}(q)$ is the model's answer and $a_{h}(q)$ is participant $h$'s answer.
The human--human (HH) analogue $\sHH[q]$ is the mean pairwise cosine similarity across all participant pairs for the same instance, serving as a per-question \emph{difficulty} measure (low $\sHH$ = high human disagreement).

\textbf{Anchoring sensitivity.} We use two complementary views of anchor dependence. First, $\Delta S \equiv S_{\vOrig} - S_{\vPron}$ measures the drop induced by pronominalization per operation or entity subgroup. Larger values indicate that responses are more tightly bound to the specific referent. % The \emph{over-degradation gap} $\Delta S_m - \Delta S_{\text{human}}$ then measures how much more (or less) a model degrades relative to humans on the same subgroup. % ; positive values indicate the model loses alignment faster than human agreement drops when the visual anchor is removed. % ( \S\ref{sec:alignment_main}; Appendix Tables~\ref{tab:op_degradation_per_model}--\ref{tab:ent_degradation_per_model};

% \textbf{Difficulty tracking.}  
Secondly, for each model $m$, we also compute Spearman $\rho$ between its per-question scores $\sHM[m,q]$ and the pooled human reference $\sHH[q]$ across all question instances. %  ($N{=}183$ 
% We use Spearman $\rho$ rather than Pearson $r$ because the relationship between $\sHM$ and $\sHH$ is monotone but not necessarily linear, and $\rho$ is robust to skewed score distributions.
A strong correlation indicates that the model's difficulty profile co-moves with the human semantic scoring on the questions. % ; a low $\pr$ indicates the model fails on structurally different questions than humans do.
% The human LOO reference is the median $\rho$ obtained by correlating each participant's leave-one-out scores against the remaining participants' pooled reference.

\textbf{Auxiliary robustness metrics.} 
Embedding-based agreement can be sensitive to answer length and response format; we therefore report a truncated robustness analysis in Appendix Figure~\ref{fig:sbert_truncation_scatter}. % the appendix (Appendix~\ref{app:model_diagnostics},  
We also report blind VQA accuracy and lexical-overlap metrics such as chrF and ROUGE-1 as supplementary diagnostics in Appendix Table~\ref{tab:comprehensive_inst_blind}.
% We therefore report lexical-overlap robustness checks (chrF, ROUGE-1) in .

\begin{figure}[!tb]
\centering
\includegraphics[width=\columnwidth]{confidence/conf_rt_sbert_triangle}
\caption{Per-question human confidence and response time (RT) vs.\ human--human SBERT ($\sHH$), shown for each control variant.
Inset: Pearson $r$ between $\sHH$ and each measure, computed separately per variant.}
\label{fig:conf_rt_sbert_triangle}
\end{figure}
% \textbf{Confidence alignment.} We also analyze token-level confidence as an
% auxiliary diagnostic of whether models express uncertainty in a human-like way.
% These analyses are not a primary selection metric in the main paper, but they
% provide a complementary view of calibration and question-level difficulty
% tracking (Appendix~\ref{app:uncertainty_appendix}).
